% BEGIN VOL08-EXPANSION VIII01
\section{Supplementary worked examples}
\label{sec:viii01-pedagogy}

\begin{example}[Straight-line homotopy in a convex target]\label{ex:viii01-ped-01}
For maps $f,g:X\to C$ into a convex subset $C\subset\mathbb R^m$, the formula $H(x,t)=(1-t)f(x)+tg(x)$ stays in $C$ and gives a homotopy.  The key point is joint continuity on $X\times I$, not merely a path for each fixed $x$.
\end{example}

\begin{example}[A rotation of the circle]\label{ex:viii01-ped-02}
For $R_\alpha(z)=e^{i\alpha}z$ on $S^1$, $H(z,t)=e^{it\alpha}z$ deforms the identity to $R_\alpha$.  Every time slice is even a homeomorphism, although homotopy theory does not require this stronger property.
\end{example}

\begin{example}[Relative homotopy of paths]\label{ex:viii01-ped-03}
If paths $\alpha,\beta:I\to Y$ have the same endpoints, a homotopy relative to $\{0,1\}$ keeps both endpoints fixed for every time.  This is exactly the equivalence relation later used to define $\pi_1(Y,y_0)$.
\end{example}

\section{Graded supplementary exercises}
The following exercises progress from direct computation and construction to proof, hypothesis testing, application, and synthesis.

\subsection*{Standard computations and constructions}

\begin{exercise}[Endpoint equations]\label{exr:viii01-ped-01}
State the two endpoint equations for $H:f\simeq g$.
\end{exercise}
\begin{hint}
Restrict $H$ to $X\times\{0\}$ and $X\times\{1\}$.
\end{hint}
\begin{solution}
They are $H(x,0)=f(x)$ and $H(x,1)=g(x)$ for every $x\in X$.
\end{solution}

\begin{exercise}[Reverse a homotopy]\label{exr:viii01-ped-02}
Given $H:f\simeq g$, write a homotopy from $g$ to $f$.
\end{exercise}
\begin{hint}
Reverse the interval parameter.
\end{hint}
\begin{solution}
Use $\overline H(x,t)=H(x,1-t)$.
\end{solution}

\begin{exercise}[Concatenate homotopies]\label{exr:viii01-ped-03}
Write the standard concatenation of $H:f\simeq g$ and $K:g\simeq h$.
\end{exercise}
\begin{hint}
Give each deformation half of $I$.
\end{hint}
\begin{solution}
Use $H(x,2t)$ for $t\le1/2$ and $K(x,2t-1)$ for $t\ge1/2$; the gluing lemma gives continuity.
\end{solution}

\begin{exercise}[Convex null-homotopy]\label{exr:viii01-ped-04}
Construct a null-homotopy of $f:X\to\mathbb R^m$.
\end{exercise}
\begin{hint}
Choose a constant target point $c$.
\end{hint}
\begin{solution}
$H(x,t)=(1-t)f(x)+tc$ contracts $f$ to the constant map $c$.
\end{solution}

\begin{exercise}[Based condition]\label{exr:viii01-ped-05}
What additional equation makes a homotopy based at $x_0\mapsto y_0$?
\end{exercise}
\begin{hint}
Keep the base point fixed for all time.
\end{hint}
\begin{solution}
Require $H(x_0,t)=y_0$ for every $t\in I$.
\end{solution}

\subsection*{Proofs}

\begin{exercise}[Equivalence relation]\label{exr:viii01-ped-06}
Prove homotopy is an equivalence relation on maps $X\to Y$.
\end{exercise}
\begin{hint}
Use constant, reversed, and concatenated homotopies.
\end{hint}
\begin{solution}
The constant homotopy gives reflexivity, $t\mapsto1-t$ gives symmetry, and piecewise concatenation plus the gluing lemma gives transitivity.
\end{solution}

\begin{exercise}[Postcomposition]\label{exr:viii01-ped-07}
Prove $f\simeq g$ implies $u\circ f\simeq u\circ g$ for continuous $u:Y\to Z$.
\end{exercise}
\begin{hint}
Compose the homotopy with $u$.
\end{hint}
\begin{solution}
If $H$ joins $f$ to $g$, then $(x,t)\mapsto u(H(x,t))$ joins the composites.
\end{solution}

\begin{exercise}[Precomposition]\label{exr:viii01-ped-08}
Prove $f\simeq g$ implies $f\circ v\simeq g\circ v$ for continuous $v:W\to X$.
\end{exercise}
\begin{hint}
Insert $v$ into the first variable.
\end{hint}
\begin{solution}
Use $(w,t)\mapsto H(v(w),t)$.
\end{solution}

\begin{exercise}[Product homotopy]\label{exr:viii01-ped-09}
Prove products of homotopic maps are homotopic.
\end{exercise}
\begin{hint}
Pair the two homotopies at the same time.
\end{hint}
\begin{solution}
If $H$ and $K$ are the homotopies, $((x,w),t)\mapsto(H(x,t),K(w,t))$ is continuous and has the required endpoints.
\end{solution}

\subsection*{Counterexamples and hypothesis tests}

\begin{exercise}[Pointwise paths suffice]\label{exr:viii01-ped-10}
Test: choosing a path from $f(x)$ to $g(x)$ for every $x$ always defines a homotopy.
\end{exercise}
\begin{hint}
Joint continuity is the missing condition.
\end{hint}
\begin{solution}
False.  The paths can vary discontinuously with $x$.
\end{solution}

\begin{exercise}[Free equals based]\label{exr:viii01-ped-11}
Test: freely homotopic based maps are always based-homotopic.
\end{exercise}
\begin{hint}
The base point may move during a free homotopy.
\end{hint}
\begin{solution}
False in general; based homotopy imposes the additional fixed-basepoint condition.
\end{solution}

\begin{exercise}[Homotopy implies isotopy]\label{exr:viii01-ped-12}
Test: every homotopy between embeddings is an isotopy.
\end{exercise}
\begin{hint}
Intermediate time slices need not remain embeddings.
\end{hint}
\begin{solution}
False.  Homotopy only asks for continuous time slices, which may fold or identify points.
\end{solution}

\subsection*{Applications and investigations}

\begin{exercise}[Deformation planning]\label{exr:viii01-ped-13}
Explain why convexity is an algorithmically useful certificate for constructing homotopies.
\end{exercise}
\begin{hint}
It makes line segments stay in the target.
\end{hint}
\begin{solution}
Convexity guarantees the straight-line interpolation remains admissible, giving an explicit continuous deformation with no obstruction search.
\end{solution}

\begin{exercise}[Cylinder viewpoint]\label{exr:viii01-ped-14}
Interpret a homotopy as an extension problem from the ends of a cylinder.
\end{exercise}
\begin{hint}
The prescribed data live on $X\times\{0,1\}$.
\end{hint}
\begin{solution}
A homotopy is a map $X\times I\to Y$ extending the two endpoint maps on the boundary copies of $X$.
\end{solution}

\subsection*{Challenge problems}

\begin{exercise}[Punctured-plane warning]\label{exr:viii01-ped-15}
Why can the straight-line homotopy from $x$ to $-x$ fail in $\mathbb R^2\setminus\{0\}$?
\end{exercise}
\begin{hint}
Evaluate at $t=1/2$.
\end{hint}
\begin{solution}
It passes through the excluded origin, so it is not a homotopy in the punctured target.
\end{solution}

\begin{exercise}[Relative synthesis]\label{exr:viii01-ped-16}
Construct a homotopy of a map on a disk while fixing its boundary, and state exactly what must be checked.
\end{exercise}
\begin{hint}
Use a formula that agrees with the original map on $S^1$ for all $t$.
\end{hint}
\begin{solution}
A valid construction must be jointly continuous on $D^2\times I$, have the prescribed endpoint maps, and satisfy $H(a,t)=f(a)$ for every $a\in S^1$ and all $t$.
\end{solution}

% END VOL08-EXPANSION VIII01
